%% titre: Quiz 6 — Nombres complexes
%% niveau: Licence 1
%% matiere: Algèbre 1
%% annee: 2026-2027
%% ordre: 6
%% description: Module, argument, formes trigonométrique et exponentielle, racines, linéarisation.

\begin{question}
Que vaut $(1+i)^2$ ?
\begin{reponses}
  \mauvaise $2$
  \bonne $2i$
  \mauvaise $-2i$
  \mauvaise $1+2i$
\end{reponses}
\begin{explication}
$(1+i)^2=1+2i+i^2=2i$.
\end{explication}
\end{question}

\begin{question}
Le module de $3-4i$ est :
\begin{reponses}
  \mauvaise $1$
  \mauvaise $7$
  \bonne $5$
  \mauvaise $25$
\end{reponses}
\begin{explication}
$|3-4i|=\sqrt{3^2+(-4)^2}=\sqrt{25}=5$.
\end{explication}
\end{question}

\begin{question}
Un argument de $-1+i$ est :
\begin{reponses}
  \mauvaise $\dfrac{\pi}{4}$
  \mauvaise $-\dfrac{\pi}{4}$
  \bonne $\dfrac{3\pi}{4}$
  \mauvaise $-\dfrac{3\pi}{4}$
\end{reponses}
\begin{explication}
Le point d'affixe $-1+i$ est dans le deuxième quadrant : $|z|=\sqrt2$, $\cos\theta=-\frac{1}{\sqrt2}$ et $\sin\theta=\frac1{\sqrt2}$, d'où $\theta=\frac{3\pi}{4}\ (2\pi)$.
\end{explication}
\end{question}

\begin{question}
Quelles affirmations sont vraies pour tous $z,z_1,z_2\in\C$ ? (plusieurs réponses possibles)
\begin{reponses}
  \bonne $z\bar z=|z|^2$
  \bonne $|z_1z_2|=|z_1|\,|z_2|$
  \mauvaise $|z_1+z_2|=|z_1|+|z_2|$
  \mauvaise $\arg(z_1+z_2)=\arg z_1+\arg z_2$
\end{reponses}
\begin{explication}
Le module est multiplicatif, mais seule l'inégalité triangulaire $|z_1+z_2|\le|z_1|+|z_2|$ est vraie en général. C'est pour le \emph{produit} que les arguments s'additionnent.
\end{explication}
\end{question}

\begin{question}
La formule de Moivre affirme que, pour $n\in\Z$ :
\begin{reponses}
  \bonne $(\cos\theta+i\sin\theta)^n=\cos(n\theta)+i\sin(n\theta)$
  \mauvaise $(\cos\theta+i\sin\theta)^n=\cos^n\theta+i\sin^n\theta$
  \mauvaise $(\cos\theta+i\sin\theta)^n=n\cos\theta+in\sin\theta$
  \mauvaise $(\cos\theta+i\sin\theta)^n=\cos(\theta^n)+i\sin(\theta^n)$
\end{reponses}
\end{question}

\begin{question}
Les racines quatrièmes de l'unité sont :
\begin{reponses}
  \mauvaise $\{1,-1\}$
  \mauvaise $\{1,i\}$
  \bonne $\{1,i,-1,-i\}$
  \mauvaise $\{1,2,3,4\}$
\end{reponses}
\begin{explication}
Ce sont les $e^{2ik\pi/4}=i^k$ pour $k=0,1,2,3$.
\end{explication}
\end{question}

\begin{question}
Pour tout réel $\theta$, $\cos\theta$ est égal à :
\begin{reponses}
  \bonne $\dfrac{e^{i\theta}+e^{-i\theta}}{2}$
  \mauvaise $\dfrac{e^{i\theta}-e^{-i\theta}}{2}$
  \mauvaise $\dfrac{e^{i\theta}-e^{-i\theta}}{2i}$
  \mauvaise $\dfrac{e^{i\theta}+e^{-i\theta}}{2i}$
\end{reponses}
\begin{explication}
Formule d'Euler. La troisième expression est $\sin\theta$.
\end{explication}
\end{question}

\begin{question}
Les solutions de $z^3=8i$ sont :
\begin{reponses}
  \mauvaise $2i$ seulement
  \bonne $\sqrt3+i,\ -\sqrt3+i,\ -2i$
  \mauvaise $2,\ 2j,\ 2j^2$
  \mauvaise $\sqrt3-i,\ -\sqrt3-i,\ 2i$
\end{reponses}
\begin{explication}
$8i=8e^{i\pi/2}$ donc $z_k=2e^{i(\pi/6+2k\pi/3)}$, $k=0,1,2$ : $2e^{i\pi/6}=\sqrt3+i$, $2e^{5i\pi/6}=-\sqrt3+i$, $2e^{3i\pi/2}=-2i$.
\end{explication}
\end{question}
